\subsection{网络架构：PMDF-Net}
\label{sec:architecture}

PMDF-Net 是一个全卷积编码器-解码器网络，在时域和时频域联合处理激光超声信号。
该架构采用双分支编码器设计配合共享解码器，在重构去噪时域波形之前实现跨域互补特征提取。

\textbf{整体结构。}PMDF-Net 包含三个阶段：时域编码器分支、时频编码器分支、瓶颈处的融合模块以及共享解码器。
两个编码器分支均逐步对输入表示进行下采样以捕获多尺度特征，然后解码器通过跳跃连接对上采样并融合信息以重构去噪输出。网络接收预处理后的输入信号$x$并输出去噪波形估计$\hat{y}$。

\textbf{时域分支。}时域分支直接在一维输入波形上运作。它由大约四到五个编码器块组成，每个块包含一个核大小为 3 或 5 的一维卷积层，随后是批归一化和 ReLU 激活。
下采样通过步长为 2 的卷积层执行，逐步减半时间分辨率同时增加通道维度，以在越来越大的空间尺度上捕获特征。
解码器以对应的上采样块数（转置卷积或步进上采样）和 U-Net 风格的跳跃连接镜像该结构到编码器特征图\cite{ronneberger2015unet}。这些跳跃连接保留了通过连续下采样过程将被丢失的高频波形细节。时域分支采用残差学习\cite{zhang2017denoising}而非直接预测干净信号：它预测噪声残差$r$，使得$\hat{y} = x - r$。这种公式通过让网络聚焦于噪声和干净信号之间的差异来简化学习目标，已被证明可以改善卷积架构的收敛速度和去噪性能。

\textbf{时频分支。}时频分支首先通过 STFT 将输入波形变换为时频表示。使用第~\ref{sec:preprocessing}节中的预处理参数，STFT 使用长度为$W = 256$的汉宁窗，跳跃大小$H = 64$，得到频率分辨率$\Delta f \approx 60$\,kHz 和时间分辨率约$0.8$\,$\mu$s 的频谱图。然后该分支将 STFT 幅度频谱图——作为单通道 2D 图像处理——送入 2D 卷积编码器-解码器。编码器应用核大小为$3 \times 3$或$5 \times 5$的 2D 卷积、批归一化和 ReLU，逐步下采样频谱图以同时捕获低频色散分量和较高频到达。
解码器通过转置卷积上采样，并在编码器处通过跳跃连接进行拼接，遵循与时域分支相同的 U-Net 拓扑。
编码器-解码器处理后，逆 STFT 重构时域分量，该分量在融合模块处与时域分支的特征进行拼接。

\textbf{融合模块。}在瓶颈处，两个分支的特征表示通过沿通道维度的拼接进行融合，得到组合特征张量，同时编码时域波形结构和时频谱特征。$1 \times 1$（逐点）卷积将拼接后的通道数压缩为紧凑表示，随后是 ReLU 非线性，使融合能够学习两个域的非平凡组合。这个融合后的瓶颈表示被传递给共享解码器。

\textbf{共享解码器。}解码器接收融合后的瓶颈表示，并通过转置卷积逐步上采样至原始信号长度。跳跃连接在每个尺度注入高分辨率特征，保留输出波形中的精细时间细节。
由于网络是全卷积的，PMDF-Net 接受任意长度的输入信号，只要长度是总下采样因子（对于四块架构为$2^4 = 16$）的倍数。

\textbf{关键设计决策。}全卷积设计确保网络无需修改即可处理变长信号，使其适用于激光超声波形的实时或批量处理。两个分支中的跳跃连接遵循 U-Net 原理，通过编码过程保留高频信息。
不存在递归层使推理快速，并避免了 RNN 去噪器中常见的梯度传播挑战。
总参数量保持在约 1--5 百万的参数范围内，足够大以学习复杂噪声模式，同时保持在桌面 GPU 推理的实际内存和计算约束内。

\textbf{输入和输出。}网络接受固定长度$L = 1024$样本的预处理一维信号$x$作为主要输入。时频分支额外需要预处理期间计算的$x$的 STFT。输出是与输入相同长度的去噪波形估计$\hat{y}$，可从中直接计算噪声残差$\hat{r} = x - \hat{y}$。